\justifying
\NSFCBodyText

\subsubsection{年度研究计划}

\textbf{第一年度：货物多模态可靠感知与实验基线。}完成完整运输任务、仓储事件、测量/可用时点和数据划分规范；组织振动、倾斜、电子封志、光感与包装图像的受控事件，复现单模态、静态融合和缺失感知基线，完成异步对齐、时变可靠性估计、运输异常识别与包装损伤定位。同步建立 2G/3G 弱网轨迹、缓存/断点/能耗记录规范和温湿度、视频、烟雾、门禁的仓储采集流程。里程碑为形成可重复的数据与评价协议，并在封存运输任务上给出 H1 的支持、失败或证据不足结论。

\textbf{第二年度：弱网边缘计算与低功耗协同。}固化货物感知模型及其校准输出，研究风险信息价值、接收端已确认状态和按时送达概率；实现本地推理深度、自适应压缩、离线缓存、断点续传、优先上传与动态工作周期的联合决策。在相同事件流、网络轨迹和初始电量上，与周期上传、阈值触发、仅信息年龄和固定压缩/工作周期策略比较。同步完成仓储单源基线与多源关系图。里程碑为在高风险漏报和 P95 告警时延约束下，报告上传成功率、传输字节、推理开销和单位任务能耗，并给出 H2 的明确判定。

\textbf{第三年度：仓储多源预警与三条主线联调。}完成温湿度缓变、烟雾/门禁突发和视频安防证据的风险演化模型，建立校准的分级预警及应急响应映射；开展来源删除、时间置乱、跨时间块和完整仓储事件外层验证。联调货物感知、弱网边缘协同和仓储预警，整理失败模式和可复现实验协议。在设备、合作和数据授权落实时开展小规模运输--仓储外部验证；条件不具备时，运输先验保持为空，但仓储四类信息的独立验证照常完成。里程碑为给出 H3 的明确判定及三条主线的端到端边界分析。

\subsubsection{预期研究结果}

预期形成货物多模态可靠性自适应感知方法、风险信息价值驱动的弱网通信--边缘计算--低功耗协同方法、仓储多源风险演化与分级预警方法，以及包含可用时点约束、完整任务/事件划分、弱网轨迹回放和失败判据的可复现实验协议；形成相应算法原型、实验代码和数据/配置说明（在数据授权允许范围内开放），发表或录用相关高水平学术论文 2--3 篇，培养研究生并为跨境物流智能感知的后续研究奠定基础。

\subsubsection{学术交流与合作计划}

围绕多模态学习、时序异常检测、边缘智能和物流物联网，每年参加相关国内外学术会议或专题研讨，汇报阶段性结果并邀请同行评议实验协议；依托新疆大学现有科研平台开展校内交叉交流。在合作单位、设备和数据合规条件明确后，与物流应用方围绕弱网轨迹、端侧能耗、运输事件和仓储多源预警开展外部验证，不预设尚未落实的国际合作任务。
